\chapter{Related Work}
\label{related-work}

Diese Arbeit beschäftigt sich mit der Fragestellung, wie die Lücke zwischen Prozessmodellierung und Prozessdurchführung im industriellen Kontext geschlossen werden kann. Um sich in dieser Thematik einen Überblick zu verschaffen, werden in diesem Kapitel ähnliche Arbeiten vorgestellt. Schließlich wird die vorliegende Arbeit in diesem Rahmen eingeordnet. Für dieses Kapitel werden die folgenden Arbeiten betrachtet:

\begin{itemize}
\item \textit{Using BPMN to Model Internet of Things Behaviour within Business Process} von Domingos, Dulce, Martins und Francisco \cite{domingos_using_2022}

\item \textit{Integration of business and production layer: Modeling and execution of vertical smart factory processes} von Schärfl und Alexandru \cite{schaerfl_integration_2016}

\item \textit{Agent-Based Business Process Orchestration for IoT} von Kampik, Timotheus, Malhi, Avleen, Främling und Kary \cite{kampik_agent-based_2019}

\item \textit{Modelling and executing IoT-enhanced business processes through BPMN and microservices} von Valderas, Pedro, Torres, Victoria, Serral und Estefía \cite{valderas_modelling_2022}
\end{itemize}

\section{Using BPMN to Model Internet of Things Behaviour within Business Process (Domingos et. al.)}

\noindent Domingos und Martins stellen in ihrer Arbeit den Ansatz vor, dass BPMN-Modelle in die plattformunabhängige Programmiersprache Callas übersetzt werden. Dies ermöglicht es, dass Teile des Geschäftsprozesses auf IoT-Hardware ausgeführt werden können. Dabei müssen technische Details der Hardware nicht berücksichtigt werden, solange die IoT-Hardware eine Callas-Laufzeitumgebung ausführen kann. Die Prozessteile, welche auf der IoT-Hardware ausgeführt werden, werden manuell aussortiert und in Callas übersetzt. Dies beschreiben die Autoren als einen Aspekt, welcher in Zukunft automatisiert werden sollte. \cite{domingos_using_2022} Während Domingos und Martins mit der Callas-Programmiersprache und der dementsprechenden Laufzeitumgebung arbeiten, wird in dieser Arbeit C++ mithilfe des ESP-IDF-Frameworks verwendet, um die Möglichkeiten des ESP32 effizient zu nutzen. Auch wird statt BPMN das EPK-Modell verwendet, aufgrund der einfachen Syntax. Beispiele dafür sind die erforderliche Abfolge von Ereignis und Funktion oder auch dass ein Start- und Endereignis zwingend notwendig sind. Jedoch wird das Prinzip der Übersetzung des Prozessmodells in ein maschinenverständliches Format von Domingos und Martins in dieser Arbeit verwendet.

\newpage

\section{Integration of business and production layer: Modeling and execution of vertical smart factory processes - Schärfl, Alexandru}

\noindent In einem ähnlichen Kontext haben sich Schärfl und Alexandru mit der Überführung von S-BPM (eine Teilmenge der BPMN-Prozessnotation) in SFC (Sequential Function Chart) beschäftigt. SFC ist dabei eine visuelle Programmiersprache gemäß des IEC-61131-3-Standards für speicherprogrammierbare Steuerungen. Hierbei zeigt Schärfl die Grenze auf, dass S-BPM nicht vollständig überführt werden kann, da Spezifika bezüglich der speicherprogrammierbaren Steuerung schon auf der Modellebene berücksichtigt werden müssen. Auch wurde nach Angaben des Autors der Aspekt der Laufzeit noch nicht beleuchtet. \cite{schaerfl_integration_2016} Schärfl und Alexandru verwenden eine Abwandlung der BPMN Prozessnotation und die Zielhardware ist eine speicherprogrammierbare Steuerung. In der vorliegenden Arbeit wird hingegen statt mit einer speicherprogrammierbaren Steuerung mit einem ESP32-Mikrocontroller gearbeitet und es werden Prozessschritte zu CoAP-Endpunkten übersetzt. Die Ansteuerung von Sensorik und Aktorik ist immer noch Aufgabe des Entwicklers. Wie schon bei Domingos und Martins erwähnt, ist das Notationsformat des Prozessmodells in der vorliegenden Arbeit EPK. Jedoch hat man von Schärfl das Prinzip der Überführung von Geschäftsprozessmodellierung zu ausführbarem Code für eingebettete Hardware, zu einem bestimmten Maß übernommen. 

\section{Agent-Based Business Process Orchestration for IoT - Kampik et. al.}

\noindent In der Arbeit von Kampik et. al. wird ein Geschäftsprozess in BPMN modelliert und mit Agenten und IoT verknüpft. Dabei kommen auch menschliche Akteure ins Spiel. Standards, die dabei verwendet werden, sind O-MI und O-DF. Dabei wurde aber nicht auf echter Hardware, sondern mit
Mock-Ups gearbeitet. Technologien, die für die Orchestrierungsebene verwendet worden sind, sind Java und Node.js.
In der Diskussion wurde die Embedded-Schicht als Thema für zukünftige Forschungsarbeiten aufgegriffen. \cite{kampik_agent-based_2019} Auch wenn in der Arbeit von Kampik et. al. im Vergleich zur vorliegenden Arbeit nicht mit echter Hardware gearbeitet wurde, wird die Verwendung von IoT-Standards übernommen. Anders als bei Kampik et. al. wird in der vorliegenden Arbeit das CoAP-Protokoll verwendet.

\section{Modelling and executing IoT-enhanced business processes through BPMN and microservices - Valderas et.al.}

\noindent In der Arbeit \textit{Modelling and executing IoT-enhanced business processes through BPMN and microservices} behandeln Valderas et.al. die Ausführung von 
IoT-gekoppelten Prozessen. Dabei werden die Prozesse in BPMN notiert und mit der Camunda Execution Engine ausgeführt, welche einen Raspberry Pi ansteuert. \cite{valderas_modelling_2022} Dies ist schon näher an dem, was diese Arbeit erreicht, nur mit dem Unterschied, dass zum einen EPK als Prozessnotation genutzt wird und die Execution Engine selbst entwickelt wird. Doch das Prinzip der Ausführung eines Geschäftsmodells auf eingebetteter Hardware von Valderas et.al. ist Hauptbestandteil dieser Arbeit, um die Ziele aus Kapitel \ref{motivation-und-zielsetzung} zu erreichen.

\noindent All diese Arbeiten verfolgen das Ziel, die Lücke zwischen der Prozessmodellierung und der Ausführung auf IoT-Hardware zu schließen. Das übergeordnete Prinzip, das Prozessmodell in eine von Maschinen ausführbare Form zu übersetzen, dient der vorliegenden Arbeit als konzeptionelle Grundlage. Das Konzept wird aber durch die Entwicklung eines eigenen Compilers mitsamt Execution-Engine und Codegenerierung realisiert, statt Plattformen wie Camunda oder Callas zu nutzen. Als konkretes Werkzeug zur Prozessmodellierung wird ARIS eingesetzt, welches noch um die Ausführung und Codegenerierung erweitert wird.

\noindent Tabelle \ref{tab:related-work-comparison} im Anhang zeigt eine Vergleichsmatrix, um die Unterschiede und Gemeinsamkeiten der verwandten Arbeiten mit der vorliegenden Arbeit übersichtlich darzustellen:
